% ─────────────────────────────────────────────────────────────────────────────
%  STT-1920 — Exercices supplémentaires : ANOVA à un facteur
%  (tailles d'échantillon égales pour chaque traitement)
% ─────────────────────────────────────────────────────────────────────────────

\def\Ponderation{}
\def\SigleNumero{STT-1920}
\def\NomCours{Méthodes statistiques}
\def\DateEvaluation{}
\def\HeureEvaluation{}
\def\Session{}
\def\NomEvaluation{Exercices supplémentaires — ANOVA à un facteur}
\def\OptionsExam{11pt}

\def\Coordonnatrice{}
\def\Coordonnateur{}
\def\Enseignant{Julien Miron}
\def\Enseignante{}
\def\SectionA{}
\def\SectionB{}
\def\SectionC{}
\def\SectionD{}
\def\NombreLignes{}
\def\EnLigne{}

\def\Directives{}

\newcommand{\affichepoints}{}     % ← AVEC points
% \renewcommand{\affichepoints}{on} % ← SANS points (non vide)

\providecommand{\Gradescope}{}

\documentclass{Evaluation}
\usepackage{multicol,graphicx,booktabs}

\printanswers

\begin{document}
\begin{questions}

% ─────────────────────────────────────────────────────────────────────────────
%  Exercice 1 — Lecture d'une table d'ANOVA (QCM, une seule bonne réponse)
% ─────────────────────────────────────────────────────────────────────────────

\question[3]
Un technicien compare l'autonomie (en heures) de casques audio de $k=4$ marques différentes.
Il teste $n=5$ casques par marque.
La table d'analyse de la variance suivante est produite par \texttt{Rcmdr} :

\begin{verbatim}
> summary(ANOVAModel)
            Df   Sum Sq   Mean Sq   F value    Pr(>F)
Marque       3   110.40    36.80    7.360    0.0024
Residuals   16    80.00     5.00
\end{verbatim}

À partir de la table d'ANOVA, quelle affirmation parmi les suivantes est \textbf{vraie}?
Une seule réponse est bonne.

\begin{checkboxes}
    \choice La durée de vie moyenne des casques pour le facteur \texttt{Marque} est égale à $36{,}80$ heures.
    \choice Il y a $5$ niveaux différents pour le facteur \texttt{Marque}.
    \choice Il y a $17$ observations en tout dans le jeu de données.
    \correctchoice Au seuil $\alpha=0{,}05$, nous pouvons affirmer que l'autonomie moyenne n'est pas la même pour toutes les marques.
    \choice Au seuil $\alpha=0{,}01$, nous ne rejetons pas $H_0$.
    \choice La somme des carrés totale vaut $5{,}00$.
\end{checkboxes}

\begin{solution}
\textbf{Analyse des choix.}
\begin{itemize}
    \item \textbf{Faux.} $36{,}80$ est le carré moyen (\texttt{Mean Sq}), pas la moyenne des observations.
    \item \textbf{Faux.} Il y a $\text{Dl}_{\text{trt}}+1=3+1=4$ niveaux (4 marques).
    \item \textbf{Faux.} $N=k\times n=4\times 5=20$ observations; les degrés de liberté totaux sont $N-1=19$.
    \item \textbf{Vrai.} La valeur-$p=0{,}0024<0{,}05=\alpha$, donc on rejette $H_0$ au seuil 5\%. On peut
    affirmer que les autonomies moyennes ne sont pas toutes égales.
    \item \textbf{Faux.} $p=0{,}0024<0{,}01$, donc on rejette aussi $H_0$ au seuil 1\%.
    \item \textbf{Faux.} $SS_\text{total}=SS_\text{trt}+SS_E=110{,}40+80{,}00=190{,}40$.
\end{itemize}
\end{solution}

\pageblanche

% ─────────────────────────────────────────────────────────────────────────────
%  Exercice 2 — Table partielle : trouver AAA, BBB, CCC (QCM multi-parties)
% ─────────────────────────────────────────────────────────────────────────────

\question
Un ingénieur compare la résistance à la traction (en MPa) de $k=5$ types d'alliages métalliques.
Il fabrique $n=6$ éprouvettes par type ($N=30$ au total).
La table d'ANOVA partielle suivante a été produite; certaines quantités ont été remplacées par
\texttt{AAA}, \texttt{BBB} et \texttt{CCC}.
Nous considérons que les postulats du modèle sont respectés.

\begin{verbatim}
ANOVA
               SS       Ddl     MS        F0       Pr(>F)
Alliage     185.40      ???    AAA       5.297     0.0037
Residuals    BBB         25   8.75
Total        CCC        ???
\end{verbatim}

\begin{parts}

    \part[2] Que vaut \texttt{AAA}?
    \begin{checkboxes}
        \choice $185{,}40$
        \choice $37{,}08$
        \correctchoice $46{,}35$
        \choice $92{,}70$
        \choice $8{,}75$
        \choice Il n'est pas possible de trouver la valeur de \texttt{AAA} avec les informations données.
    \end{checkboxes}

    \begin{solution}
    $\text{Dl}_\text{trt}=k-1=5-1=4$.
    \[
        \texttt{AAA}=MS_\text{trt}=\frac{SS_\text{trt}}{\text{Dl}_\text{trt}}=\frac{185{,}40}{4}=46{,}35.
    \]
    \end{solution}

    \vspace{0.5cm}

    \part[2] Que vaut \texttt{BBB}?
    \begin{checkboxes}
        \choice $25$
        \choice $8{,}75$
        \choice $185{,}40$
        \correctchoice $218{,}75$
        \choice $404{,}15$
        \choice Il n'est pas possible de trouver la valeur de \texttt{BBB} avec les informations données.
    \end{checkboxes}

    \begin{solution}
    \[
        \texttt{BBB}=SS_E=MS_E\times \text{Dl}_E=8{,}75\times 25=218{,}75.
    \]
    \end{solution}

    \vspace{0.5cm}

    \part[2] Combien y a-t-il d'observations en tout dans le jeu de données?
    \begin{checkboxes}
        \choice $25$
        \choice $29$
        \correctchoice $30$
        \choice $31$
        \choice Il n'est pas possible de répondre avec les informations données.
    \end{checkboxes}

    \begin{solution}
    $\text{Dl}_E=N-k=25$, donc $N=25+k=25+5=30$.
    On vérifie : $k\times n = 5\times 6=30$. $\checkmark$
    \end{solution}

    \vspace{0.5cm}

    \part[2] Au seuil de 5\%, quelle conclusion pouvons-nous tirer de la table d'ANOVA?
    Une seule réponse est bonne.
    \begin{checkboxes}
        \choice Il n'est pas possible de tirer une conclusion car certaines valeurs sont remplacées par \texttt{???}.
        \choice Toutes les résistances moyennes sont significativement égales.
        \choice Toutes les résistances moyennes sont significativement différentes.
        \correctchoice La résistance moyenne d'au moins un type d'alliage est significativement différente de celle d'au moins un autre type.
        \choice Il faudrait davantage d'éprouvettes pour pouvoir conclure.
    \end{checkboxes}

    \begin{solution}
    La valeur-$p=0{,}0037 < 0{,}05=\alpha$, donc on \textbf{rejette $H_0$}.
    On conclut que la résistance moyenne d'au moins un type d'alliage est significativement différente de celle d'au moins un autre type. (L'ANOVA ne précise pas lesquels.)
    \end{solution}

\end{parts}

\pageblanche

% ─────────────────────────────────────────────────────────────────────────────
%  Exercice 3 — Calculs à partir des statistiques de groupe
% ─────────────────────────────────────────────────────────────────────────────

\question
Un biologiste étudie l'effet de l'intensité lumineuse sur la durée de germination (en jours) de semences.
Il soumet $n=8$ semences à chacune des $k=3$ conditions d'éclairage ($N=24$ au total) et obtient les résultats résumés ci-dessous.
On suppose que les postulats de l'ANOVA à un facteur sont satisfaits.

\medskip
\begin{center}
\renewcommand{\arraystretch}{1.4}
\begin{tabular}{|l|c|c|}
\hline
\textbf{Condition} & \textbf{Moyenne} $\bar{y}_i$ & \textbf{Variance} $s_i^2$ \\
\hline
Faible & $12{,}0$ & $4{,}0$ \\
Modérée & $15{,}0$ & $5{,}0$ \\
Forte & $18{,}0$ & $3{,}0$ \\
\hline
\end{tabular}
\end{center}
\medskip

\begin{parts}

    \part[2] Calculez $SS_\text{trt}$, la somme des carrés des traitements.
    \textbf{\textit{Montrez votre démarche.}}

    \begin{solution}
    La moyenne globale est $\bar{y}_{..}=\dfrac{12{,}0+15{,}0+18{,}0}{3}=15{,}0$.
    \[
        SS_\text{trt}=n\sum_{i=1}^{k}\bigl(\bar{y}_i-\bar{y}_{..}\bigr)^2
        =8\bigl[(12{,}0-15)^2+(15{,}0-15)^2+(18{,}0-15)^2\bigr]
        =8\bigl[9+0+9\bigr]=144.
    \]
    \end{solution}

    \vspace{6cm}

    \part[2] Calculez $MS_E$, le carré moyen des résidus.
    \textbf{\textit{Montrez votre démarche.}}

    \begin{solution}
    \[
        SS_E=\sum_{i=1}^{k}(n-1)\,s_i^2=(8-1)(4{,}0+5{,}0+3{,}0)=7\times 12=84.
    \]
    \[
        MS_E=\frac{SS_E}{N-k}=\frac{84}{24-3}=\frac{84}{21}=4{,}0.
    \]
    \end{solution}

    \vspace{6cm}

    \part[3] Calculez la statistique de test $F_0$.
    \textbf{\textit{Montrez votre démarche.}}

    \begin{solution}
    \[
        MS_\text{trt}=\frac{SS_\text{trt}}{k-1}=\frac{144}{2}=72{,}0.
    \]
    \[
        F_0=\frac{MS_\text{trt}}{MS_E}=\frac{72{,}0}{4{,}0}=18{,}0.
    \]
    \end{solution}

    \vspace{6cm}

    \part[3] En utilisant la valeur critique $F_{0{,}05\,;\,2\,;\,21}=3{,}467$, quelle est la conclusion
    du test au seuil $\alpha=0{,}05$? Rédigez votre conclusion en lien avec le contexte.

    \begin{solution}
    Comme $F_0=18{,}0 > 3{,}467 = F_{0{,}05\,;\,2\,;\,21}$, on \textbf{rejette $H_0$} au seuil 5\%.
    Il y a suffisamment de preuves pour affirmer que la durée de germination moyenne n'est pas la même pour toutes les conditions d'éclairage.
    \end{solution}

    \vspace{4cm}

\end{parts}

\pageblanche

% ─────────────────────────────────────────────────────────────────────────────
%  Exercice 4 — Vrai ou Faux sur les concepts de l'ANOVA
% ─────────────────────────────────────────────────────────────────────────────

\question
Un pharmacien compare la réduction de la douleur (en points sur une échelle de 10) de $k=4$
médicaments analgésiques. Chaque médicament est administré à $n=10$ patients ($N=40$ au total).
La table ci-dessous est obtenue :

\begin{verbatim}
> summary(ANOVAModel)
               Df    Sum Sq    Mean Sq    F value    Pr(>F)
Medicament      3    241.50     80.50      2.483    0.0763
Residuals      36   1167.48     32.43
\end{verbatim}

Pour chaque énoncé qui suit, indiquez si l'énoncé est \textbf{vrai} ou \textbf{faux}.

\medskip

\begin{parts}

    \part[2] L'hypothèse nulle du test est $H_0: \mu_1=\mu_2=\mu_3=\mu_4$, où $\mu_i$ désigne la réduction
    moyenne de la douleur pour le médicament $i$.

    \medskip
    \begin{oneparcheckboxes}
        \correctchoice Vrai
        \choice Faux
    \end{oneparcheckboxes}
    \medskip

    \begin{solution}
    \textbf{Vrai.} L'ANOVA à un facteur teste l'égalité de toutes les moyennes de groupe.
    \end{solution}

    \part[2] Au seuil $\alpha=0{,}05$, on rejette $H_0$.

    \medskip
    \begin{oneparcheckboxes}
        \choice Vrai
        \correctchoice Faux
    \end{oneparcheckboxes}
    \medskip

    \begin{solution}
    \textbf{Faux.} La valeur-$p = 0{,}0763 > 0{,}05 = \alpha$, donc on \textbf{ne rejette pas} $H_0$ au seuil 5\%.
    \end{solution}

    \part[2] Au seuil $\alpha=0{,}10$, on rejette $H_0$.

    \medskip
    \begin{oneparcheckboxes}
        \correctchoice Vrai
        \choice Faux
    \end{oneparcheckboxes}
    \medskip

    \begin{solution}
    \textbf{Vrai.} La valeur-$p = 0{,}0763 < 0{,}10 = \alpha$.
    \end{solution}

    \part[2] Le nombre de degrés de liberté totaux est $39$.

    \medskip
    \begin{oneparcheckboxes}
        \correctchoice Vrai
        \choice Faux
    \end{oneparcheckboxes}
    \medskip

    \begin{solution}
    \textbf{Vrai.} $\text{Dl}_\text{total}=N-1=40-1=39$.
    On vérifie : $\text{Dl}_\text{trt}+\text{Dl}_E=3+36=39$. $\checkmark$
    \end{solution}

\end{parts}

\pageblanche

% ─────────────────────────────────────────────────────────────────────────────
%  Exercice 5 — Calculs à partir de SST et SS_E
% ─────────────────────────────────────────────────────────────────────────────

\question
Un ingénieur alimentaire compare la teneur en sucres (en g/100 mL) de boissons gazeuses
produites par $k=6$ fabricants différents.
Il sélectionne $n=5$ bouteilles par fabricant ($N=30$ au total).
Il obtient les sommes des carrés suivantes :
\[
    SS_\text{total}=890{,}40, \qquad SS_E=540{,}00.
\]
On suppose que les postulats du modèle sont satisfaits.

\begin{parts}

    \part[2] Écrivez les hypothèses du test ANOVA correspondant à cette situation.

    \begin{solution}
    Soit $\mu_i$ la teneur moyenne en sucres des boissons du fabricant $i$, pour $i=1,\ldots,6$.
    \[
        H_0: \mu_1=\mu_2=\mu_3=\mu_4=\mu_5=\mu_6
        \quad\text{contre}\quad
        H_1: \text{au moins deux moyennes sont différentes}.
    \]
    \end{solution}

    \vspace{3cm}

    \part[2] Calculez $SS_\text{trt}$, la somme des carrés des traitements.
    \textbf{\textit{Montrez votre démarche.}}

    \begin{solution}
    \[
        SS_\text{trt}=SS_\text{total}-SS_E=890{,}40-540{,}00=350{,}40.
    \]
    \end{solution}

    \vspace{3cm}

    \part[2] Calculez $MS_\text{trt}$ et $MS_E$.
    \textbf{\textit{Montrez votre démarche.}}

    \begin{solution}
    \[
        MS_\text{trt}=\frac{SS_\text{trt}}{k-1}=\frac{350{,}40}{6-1}=\frac{350{,}40}{5}=70{,}08.
    \]
    \[
        MS_E=\frac{SS_E}{N-k}=\frac{540{,}00}{30-6}=\frac{540{,}00}{24}=22{,}50.
    \]
    \end{solution}

    \vspace{3cm}

    \part[2] Calculez la statistique de test $F_0$.
    \textbf{\textit{Montrez votre démarche.}}

    \begin{solution}
    \[
        F_0=\frac{MS_\text{trt}}{MS_E}=\frac{70{,}08}{22{,}50}=3{,}114.
    \]
    \end{solution}

    \vspace{3cm}

    \part[2] En utilisant la valeur critique $F_{0{,}05\,;\,5\,;\,24}=2{,}621$, quelle est la conclusion
    du test au seuil $\alpha=0{,}05$? Rédigez votre conclusion en lien avec le contexte.

    \begin{solution}
    Comme $F_0=3{,}114 > 2{,}621=F_{0{,}05\,;\,5\,;\,24}$, on \textbf{rejette $H_0$} au seuil 5\%.
    Il y a suffisamment de preuves pour affirmer que la teneur moyenne en sucres n'est pas la même pour tous les fabricants.
    \end{solution}

    \vspace{4cm}

\end{parts}

\end{questions}
\end{document}
